\section{Introduction}
\label{sec:intro}

Agentic LLM applications interleave generation with branching.
A coding agent proposes several candidate edits or tool calls for the same repository context, executes one, and discards the rest~\cite{react,sweagent,openhands}; a reasoning harness expands a Tree-of-Thoughts frontier or a self-consistency pool from one partial solution~\cite{tot,selfconsistency,cot}.
In both cases the serving system observes the same pattern at every step: a committed context (the \emph{trunk}) is extended by $k$ candidate continuations (the \emph{residuals}), and at most one of them (the \emph{spine}) survives into the next step.
Most of the candidates are discarded~\cite{gsm8k,math,esc}, so the efficiency of an agentic serving system depends on how early, and at which stage, a non-surviving candidate stops consuming resources.

Production stacks increasingly serve such workloads with prefill--decode (PD) disaggregation, which places prompt processing and token generation on separate instances and moves the KV cache between them through a connector~\cite{distserve,splitwise}.
Decode-time path pruning has developed in parallel: ESC stops sampling once a window agrees~\cite{esc}, Speculative Rejection halts low-reward partial generations at a horizon~\cite{specrej}, DPTS early-stops children after a mini-step~\cite{dpts}, and STOP learns a token that predicts futility from a short prefix~\cite{stoploss}.
Combined naively, these two lines produce a structural mismatch.
A pruning decision is taken on the decode instance, because every existing signal reads generated tokens, while the cost the decision is meant to avoid---prefill compute on the prefill instance and KV transfer across the connector---has already been incurred.
Pruning is therefore \emph{stage-local}: each stage can only discard work that has already reached it.

Prefix caching does not remove this mismatch.
Automatic prefix caching (APC) in vLLM and RadixAttention in SGLang reuse the trunk whenever its blocks are already published and resident~\cite{vllm,sglang}; whether $k$ siblings issued together share the trunk depends on the cache regime and on the engine's scheduling (\S\ref{sec:regimes}).
Even when the trunk is fully shared, prefix caching neither rejects a residual nor controls which pages cross the connector.
Program-level schedulers and KV-retention systems for agents~\cite{autellix,continuum,mori,conserve} manage the context across tool calls, and ForkKV shares a base cache across LoRA agents~\cite{forkkv}; none of them decides, per candidate, whether it should be prefilled, transferred, or decoded.

\begin{figure*}[t]
\centering
\resizebox{0.78\textwidth}{!}{\begin{tikzpicture}[
  x=1cm, y=1cm,
  font=\sffamily,
  line cap=round, line join=round,
]
\definecolor{dsctrl}{RGB}{233,150,170}
\definecolor{dsmodel}{RGB}{245,215,95}
\definecolor{dsrun}{RGB}{170,220,230}
\definecolor{dsgpu}{RGB}{95,185,115}
\definecolor{dsinst}{RGB}{235,235,235}
\definecolor{dsedge}{RGB}{55,55,55}
\definecolor{dspre}{RGB}{255,200,140}
\definecolor{dsdec}{RGB}{180,210,255}

\path (0,0.75) rectangle (16.2,9.0);

\foreach \x in {1.35, 2.15, 2.95, 3.75, 4.55} {
  \draw[-{Latex[length=2.4mm,width=1.9mm]}, line width=1.0pt, draw=dsedge]
    (\x,8.75) -- (\x,8.15);
}
\node[font=\footnotesize\sffamily] at (2.95,8.92) {Requests};

\draw[draw=dsedge, fill=dsctrl, line width=0.8pt]
  (0.70,7.35) rectangle (15.50,8.05);
\node[font=\large\sffamily\bfseries] at (8.10,7.82) {Controller};
\node[font=\scriptsize\sffamily] at (8.10,7.50) {decide before any kernel: residency $m_i^{\Pin},m_i^{\Din}$ and score $s_i$ $\rightarrow$ $a_i\in\{\Hit,\Admit,\Probe,\Reject\}$};

\draw[-{Latex[length=2.4mm,width=1.9mm]}, line width=1.05pt, draw=dsedge]
  (3.15,7.35) -- (3.15,6.85);
\draw[-{Latex[length=2.4mm,width=1.9mm]}, line width=1.05pt, draw=dsedge]
  (13.05,7.35) -- (13.05,6.85);

\draw[draw=dsedge, fill=dspre, line width=0.75pt]
  (1.15,6.05) rectangle (5.15,6.75);
\node[font=\small\sffamily\bfseries] at (3.15,6.40) {Prefill Pruner};

\draw[draw=dsedge, fill=dsdec, line width=0.75pt]
  (11.05,6.05) rectangle (15.05,6.75);
\node[font=\small\sffamily\bfseries] at (13.05,6.40) {Decoding Pruner};

\draw[-{Latex[length=2.4mm,width=1.9mm]}, line width=1.05pt, draw=dsedge]
  (3.15,6.05) -- (3.15,5.55);
\draw[-{Latex[length=2.4mm,width=1.9mm]}, line width=1.05pt, draw=dsedge]
  (13.05,6.05) -- (13.05,5.55);

\draw[draw=dsedge, fill=dsinst, line width=0.7pt]
  (1.15,0.85) rectangle (6.15,5.05);
\draw[draw=dsedge, fill=dsinst, line width=0.7pt]
  (0.90,1.10) rectangle (5.90,5.30);
\draw[draw=dsedge, fill=dsinst, line width=0.85pt]
  (0.65,1.35) rectangle (5.65,5.55);

\node[font=\normalsize\sffamily\bfseries] at (3.15,5.25) {Prefill Instance};

\draw[draw=dsedge, fill=dsmodel, line width=0.7pt]
  (1.05,4.35) rectangle (5.25,4.95);
\node[font=\small\sffamily] at (3.15,4.65) {LLM Model};

\draw[draw=dsedge, fill=dsrun, line width=0.7pt]
  (1.05,1.65) rectangle (5.25,4.15);
\node[font=\small\sffamily] at (3.15,3.85) {Parallel Runtime};

\foreach \gx/\gy in {1.35/2.85, 3.25/2.85, 1.35/1.90, 3.25/1.90} {
  \draw[draw=dsedge, fill=dsgpu, line width=0.6pt]
    (\gx,\gy) rectangle ({\gx+1.50},{\gy+0.72});
  \node[font=\small\sffamily] at ({\gx+0.75},{\gy+0.36}) {GPU};
}

\draw[draw=dsedge, fill=dsinst, line width=0.7pt]
  (11.05,0.85) rectangle (16.05,5.05);
\draw[draw=dsedge, fill=dsinst, line width=0.7pt]
  (10.80,1.10) rectangle (15.80,5.30);
\draw[draw=dsedge, fill=dsinst, line width=0.85pt]
  (10.55,1.35) rectangle (15.55,5.55);

\node[font=\normalsize\sffamily\bfseries] at (13.05,5.25) {Decoding Instance};

\draw[draw=dsedge, fill=dsmodel, line width=0.7pt]
  (10.95,4.35) rectangle (15.15,4.95);
\node[font=\small\sffamily] at (13.05,4.65) {LLM Model};

\draw[draw=dsedge, fill=dsrun, line width=0.7pt]
  (10.95,1.65) rectangle (15.15,4.15);
\node[font=\small\sffamily] at (13.05,3.85) {Parallel Runtime};

\foreach \gx/\gy in {11.25/2.85, 13.15/2.85, 11.25/1.90, 13.15/1.90} {
  \draw[draw=dsedge, fill=dsgpu, line width=0.6pt]
    (\gx,\gy) rectangle ({\gx+1.50},{\gy+0.72});
  \node[font=\small\sffamily] at ({\gx+0.75},{\gy+0.36}) {GPU};
}

\fill[white]
  (5.95,2.85) -- (8.35,2.85) -- (8.35,2.45) -- (9.85,3.20) --
  (8.35,3.95) -- (8.35,3.55) -- (5.95,3.55) -- cycle;
\draw[draw=dsedge, line width=1.25pt]
  (5.95,2.85) -- (8.35,2.85) -- (8.35,2.45) -- (9.85,3.20) --
  (8.35,3.95) -- (8.35,3.55) -- (5.95,3.55) -- cycle;

\node[font=\footnotesize\sffamily\bfseries] at (7.40,4.20) {KV Cache Transfer};
\node[font=\scriptsize\sffamily, align=center] at (7.40,2.05)
  {decoder-missing pages\\[-1pt]of admitted residuals only};

\end{tikzpicture}
}
\caption{\textbf{\sys{} runtime.}
The controller reads the residency of each residual on the prefill instance ($m_i^{\Pin}$) and the decode instance ($m_i^{\Din}$) together with a model-free score $s_i$, and assigns an action $a_i\in\{\Hit,\Admit,\Probe,\Reject\}$ before any kernel runs.
The prefill instance computes only the non-resident suffix of admitted and probed residuals; the connector carries only the pages the decode instance lacks.
The decode-side pruner extends a probed residual beyond $t_0$ tokens only if its generated prefix clears $\tau_{\mathrm{dec}}$.
After commit, only the spine $x_\star$ remains resident.}
\label{fig:system}
\end{figure*}

\paragraph{PD pruning disaggregation.}
We argue that, in disaggregated agentic serving, pruning should itself be disaggregated: the \emph{decision}, its \emph{execution}, and the \emph{management of pruned state} should be placed at the stages where each is cheapest and most informative.
\sys instantiates this principle (\Cref{fig:system}).
\begin{itemize}[leftmargin=*,itemsep=1pt,topsep=2pt]
\item \emph{Decision placement.} A controller-side admission policy assigns each residual one of four actions before any kernel runs, using only signals available at that point: the page-aligned residency of the residual on both instances and a model-free score of the residual text (\S\ref{sec:app}).
\item \emph{Cross-stage coordination.} Residuals that cannot be resolved before prefill are \emph{probed}: they are prefilled, decoded for $t_0$ tokens, and extended only if the generated prefix clears a decode-side threshold $\tau_{\mathrm{dec}}$ that is calibrated separately from the prefill-side threshold $\tau_{\mathrm{pre}}$ (\S\ref{sec:split}). Decode-time pruners such as ESC, Speculative Rejection, and DPTS operate unchanged on the admitted set.
\item \emph{Pruning-aware KV management.} A copy-on-write context tree shares the trunk across siblings, transfers only decoder-missing pages, frees residual pages on abort, and retains the committed spine as the trunk of the next agent step (\S\ref{sec:state}).
\end{itemize}

\paragraph{Results.}
On Qwen3-8B GSM8K with $k{=}16$, admission in front of ESC, Speculative Rejection, or DPTS reduces computed prefill from $5008$ to $3848$ tokens ($-23\%$) for every host; committed-answer accuracy changes by $0$, $-2$, and $0$ of $16$ problems, and ESC wall-clock falls from $17.05$\,s to $14.13$\,s (\S\ref{sec:eval}).
On MATH-500, rejecting every low-scoring residual before prefill lowers any-finished coverage from $20/32$ to $13/32$, whereas a $128$-token probe restores $20/32$ at $72.2$\,s instead of $104.1$\,s.
On the full GSM8K test set with Qwen3-14B, the retained spine requires $0.72\times$ the median peak KV of APC at matched accuracy, and the ratio stays within $0.72$--$0.75\times$ across six model families.
All measurements run on physically disaggregated vLLM prefill and decode instances connected by a KV connector.

\paragraph{Contributions.}
\begin{itemize}[leftmargin=*,itemsep=1pt,topsep=2pt]
\item We formulate pruning in disaggregated agentic serving as a per-residual, cross-stage decision problem, and make explicit the cache regimes and residency conditions under which prefix caching does or does not share the trunk (\S\ref{sec:problem}).
\item We design \sys, which disaggregates the pruning decision (pre-kernel admission), its execution (prefill, connector, and decode-side probe), and pruned-state management (copy-on-write context tree), and we give a work and transfer accounting that separates trunk sharing, admission, and probe overhead (\S\ref{sec:design}).
\item We show on Qwen3 models that admission composes with three decode-time pruners, that a decode-side probe recovers the coverage lost by blanket pre-kernel rejection, and that spine retention lowers peak KV, while reporting committed-answer accuracy throughout (\S\ref{sec:eval}).
\end{itemize}
